\section*{8. Über ganze rationale Darstellung der Invarianten eines Systems von beliebig vielen Grundformen}
\begin{center}
Math. Ann. 77 (1916), S. 93--102
\end{center}

\begin{center}
{\Large Über ganze rationale Darstellung der Invarianten eines Systems\\
von beliebig vielen Grundformen.}\par
\vspace{1.2em}
Von\par
\vspace{0.8em}
{\scshape Emmy Noether} in Erlangen.
\end{center}

Am Schluß seines Beitrags zur Schwarz-Festschrift ``Über die Invarianten eines Systems von beliebig vielen Grundformen'' spricht D. Hilbert die Vermutung aus, daß sich \emph{diese Invarianten ganz und rational durch die endlich vielen Invarianten des Systems $(J,PJ)$ darstellen lassen}, wo $J$ das volle Invariantensystem von den linear unabhängigen der Grundformen bedeutet, und die Invarianten $PJ$ aus $J$ durch Polarprozesse hervorgehen.

Ich zeige im folgenden in I, daß diese Vermutung tatsächlich zutrifft, und zwar als einfache Folge des Reduktionssatzes, \emph{daß jede Form von beliebig vielen Reihen von je $n$ Variabeln sich darstellen läßt als Summe von Polaren von Formen, die nur $n$ feste Reihen enthalten und ihrerseits aus der Ausgangsform durch Polarprozesse abgeleitet sind}. Dieser Reduktionssatz ist ein zuerst von F. Mertens formulierter Spezialfall der von A. Capelli, F. Mertens und J. Deruyts gegebenen Verallgemeinerung der Clebsch-Gordanschen Reihenentwicklung auf Formen von beliebig vielen Reihen von je $n$ Variabeln; eine Verallgemeinerung, die durch formale Umformung der Differentiationsprozesse bewiesen ist.\NoetherSrcNote{*)}{F. Mertens: ``Über eine Formel der Determinantentheorie.'' (Formel 5) Sitzb. d. Ak. d. Wiss. Wien, Bd. 91, II. Abt. (1885), und ausführlicher (mit Anwendungen) in: ``Über ganze Funktionen von $m$ Systemen von je $n$ Unbestimmten''. Monatsh. f. Math. u. Phys. 4. Jahrg. (1893). Für einige der Capellischen Beweise (seit 1882) vgl. etwa: Math. Ann. 37, oder die (autographierten) ``Lezioni sulla teoria delle forme algebriche'' (1902). J. Deruyts: \emph{Essai d'une theorie générale des formes algébriques}, Mém. d. l. Soc. Roy. des Sciences de Liège (2) 17 (1892).}

Ich gebe noch in II. einen mehr begrifflichen Beweis des Reduktionssatzes, indem ich die Frage als ein Problem der Äquivalenz von linearen Formenscharen auffasse. Diese Auffassung und ebenso einen wesentlichen Teil der benutzten Überlegungen entnehme ich einer Arbeit ``Über algebraische Modulsysteme''\NoetherSrcNote{*)}{E. Fischer: Über algebraische Modulsysteme und lineare homogene partielle Differentialgleichungen mit konstanten Koeffizienten, \S\ 1--4, J. f. M. 140 (1911).} und daran anschließenden unveröffentlichten Sätzen von E. Fischer, deren Verwendung mir dieser freundlichst gestattete.

Schließlich zeige ich noch in III., wie die entsprechenden Überlegungen auch zu den allgemeinsten Reihenentwicklungen führen.

\subsection*{I.}
Es sei $\theta$ eine in jeder der $N>n$ Reihen
\[
  A_1,A_2,\ldots,A_N
\]
homogene Form, wobei allgemein die Reihe $A_k$ aus den $n$ Elementen
\[
  A_k^{(1)},A_k^{(2)},\ldots,A_k^{(n)}
\]
bestehen möge. Die Koeffizienten von $\theta$ können Unbestimmte oder Größen eines gegebenen Rationalitätsbereichs sein. Es bezeichne ferner $P$ eine ganze rationale Funktion --- \emph{mit rationalen Zahlkoeffizienten} --- der Polarprozesse:
\[
  P_{hk}=\left(A_h\frac{\partial}{\partial A_k}\right)
  =A_h^{(1)}\frac{\partial}{\partial A_k^{(1)}}+
   A_h^{(2)}\frac{\partial}{\partial A_k^{(2)}}+
   \cdots+
   A_h^{(n)}\frac{\partial}{\partial A_k^{(n)}},
\]
wobei unter dem Produkt $P_{hk}P_{ij}\theta$ die Anwendung von $P_{hk}$ auf $P_{ij}\theta$ zu verstehen ist.

Dann ist der erwähnte \emph{Reduktionssatz} gegeben durch die Identität
\begin{equation}
  \theta=\sum PZ,
\tag{1}
\end{equation}
wo die Formen $Z$ nur noch die Reihen
\[
  A_1,A_2,\ldots,A_n
\]
enthalten, und aus $\theta$ durch Polarprozesse $P$ abgeleitet sind.

Wir betrachten nun das Grundformensystem $(F')$, bestehend aus den $N$ Formen gleicher Ordnung und gleicher Variabelnzahl
\[
  F_1,F_2,\ldots,F_N,
\]
deren Koeffizienten bzw. als die $N$ Reihen
\[
  A_1,A_2,\ldots,A_N
\]
genommen werden. Mit $(J)$ sei das --- aus endlich vielen Invarianten bestehende --- volle Systeme von $F_1,\ldots,F_n$ bezeichnet, derart, daß jede Invariante von $F_1,\ldots,F_n$ sich ganz und rational durch die Invarianten $(J)$ ausdrücken läßt. \emph{Die zu beweisende Hilbertsche Vermutung besteht dann darin, daß für die Simultaninvarianten $S$ von $F_1,\ldots,F_N$ ein volles System durch die endlich vielen Invarianten $(J,PJ)$ gegeben ist, wo $PJ$ alle durch Polarprozesse $P$ aus $J$ ableitbaren Invarianten bedeutet.}

In der Tat wird jede rationalzahlige Simultaninvariante $S$ von $(F')$ eine in jeder der $N$ Reihen $A_1,\ldots,A_N$ homogene Form --- mit rationalen Zahlenkoeffizienten ---, auf die sich somit die Identität (1) anwenden läßt. Da die Anwendung der Polarprozesse $P$ auf $S$ bekanntlich wieder Invarianten erzeugt, werden auch die Formen $Z$ wieder Invarianten, und zwar Simultaninvarianten von $F_1,
\ldots,F_n$, da sie nur noch die Reihen $A_1,
\ldots,A_n$ enthalten; sie sind also nach Voraussetzung ganze rationale Funktionen der Invarianten $(J)$. Somit geht die Identität (1) über in
\[
  S=\sum PY,
\]
wo die $Y$ Potenzprodukte der Invarianten $(J)$ bedeuten. Die wirkliche Ausführung der Prozesse $P$ auf $Y$ besteht aber nach der Definition von $P$ in der sukzessiven, endlich oft wiederholten Ausführung der einfachen Polaroperationen $P_{hk}$, die nach der Differentiationsregel für Produkte auf Summen von Produkten aus Invarianten $(J)$ und $(PJ)$ führen. Das gleiche gilt für die Anwendung von $P_{hk}$ auf ein Produkt aus $(J)$ und $(PJ)$, und somit liefert auch $PY$ nur ganze rationale Verbindungen von $(J,PJ)$. \emph{Die ganze rationale Darstellbarkeit aller Invarianten $S$ durch $(J,PJ)$ ist damit bewiesen.}

